% Niveaux d'énergie du néon mis en jeu dans le laser hélium-néon (exercice 15, chapitre 25)
\begin{tikzpicture}[x=1cm,y=1cm]
\draw[otMuted,line width=0.4pt,-{Stealth[length=1.5mm]}] (0,0) -- (0,3.3) node[pcLegende,anchor=south] {$E$ (eV)};
\foreach \E/\y/\lab in {0/0.2/{0 (état fondamental)},16.67/1.3/{16,67},18.70/2.1/{18,70},20.66/2.9/{20,66}}{
  \draw[otNavy,line width=1pt] (0.6,\y) -- (3.2,\y);
  \node[pcLegende,anchor=west] at (3.3,\y) {\lab};}
\draw[otRed!80!black,line width=1pt,-{Stealth[length=1.8mm]},decorate,
  decoration={snake,amplitude=0.5mm,segment length=2mm,post length=2mm}] (1.9,2.88) -- (1.9,2.12);
\node[pcLegende,anchor=east,text=otRed!80!black] at (1.8,2.5) {émission laser};
\draw[otMuted,line width=0.6pt,dash pattern=on 2pt off 1.5pt] (0.6,0.6) -- (0.6,0.6);
\node[pcLegende,anchor=north west,text=otMuted] at (0.6,1.15) {(échelle non respectée)};
\end{tikzpicture}
\par\smallskip\legende{Quelques niveaux d'énergie de l'atome de néon.}
